Des de la platja, observem que la distància entre les crestes de dues onades consecutives és de 4~m. Per altra banda, tot observant una de les boies que limita la zona de bany, comptem que oscil·la 30 vegades en un minut i que el seu desplaçament vertical total, des de la posició més baixa fins al punt més alt, és de 40~cm.

\begin{apartat}{1,25}
Determineu la freqüència, la longitud d'ona i la velocitat de propagació de les ones. Escriviu l'equació que descriu el moviment de la boia en funció del temps.
\end{apartat}

\begin{apartat}{1,25}
Deduïu l'expressió de la velocitat i l'acceleració de la boia i calculeu els valors màxims de la velocitat i l'acceleració. Representeu a la quadrícula de sota l'evolució de la velocitat respecte al temps durant un període.
\end{apartat}

\figura[0.5]{fig1.png}

\nota{Per determinar la fase inicial considereu que al principi la boia es troba en la posició més alta.}
